% ======================================================================
%  Титульный лист и «Как пользоваться книгой».
%  Перед \input задать \def\fmSubtitle{...} (что входит в данный файл).
% ======================================================================
\begin{titlepage}
\pdfbookmark[0]{Титульный лист}{title}
\centering
\vspace*{18mm}
{\sffamily\large\color{fmGray} Переработанное учебное изложение\par}
\vspace{14mm}
{\sffamily\bfseries\fontsize{30}{36}\selectfont\color{fmHead}
  Физическое\\[2mm] материаловедение\par}
\vspace{6mm}
{\color{fmBlue}\rule{0.55\textwidth}{1.2pt}\par}
\vspace{6mm}
{\sffamily\LARGE\color{fmHead} Том 1. Физика твёрдого тела\par}
\vspace{10mm}
{\sffamily\Large \fmSubtitle\par}
\vfill
\begin{minipage}{0.86\textwidth}\small\raggedright
Изложение построено на материале учебника:
\emph{Физическое материаловедение: Учебник для вузов: в 6~т.} / Под общ.\ ред.\ Б.\,А.~Калина.
Т.~1. Физика твёрдого тела / Г.\,Н.~Елманов, А.\,Г.~Залужный, В.\,И.~Скрытный,
Е.\,А.~Смирнов, В.\,Н.~Яльцев. — М.: МИФИ, 2007. — 636~с.

\medskip
Текст переписан заново: материал разбит на короткие смысловые шаги,
выводы расписаны подробно, добавлены рисунки, числовые примеры,
итоги разделов и вопросы для самопроверки. Нумерация глав и разделов
совпадает с исходным учебником, поэтому их удобно читать параллельно.
Нумерация рисунков, формул и таблиц~— своя.
\end{minipage}
\vspace{12mm}
\end{titlepage}

\cleardoublepage
\chapter*{Как пользоваться этой книгой}
\addcontentsline{toc}{chapter}{Как пользоваться этой книгой}
\markboth{Как пользоваться этой книгой}{Как пользоваться этой книгой}

Основной текст читается подряд. Цветные блоки помогают понять, \emph{какую роль}
играет фрагмент: это новое понятие, ключевой вывод, расчёт или
дополнительный материал. Ниже~— их список.

\begin{orientir}
Такой рамкой открывается каждый раздел: что в нём будет и зачем это нужно.
\end{orientir}

\begin{opredelenie}
Новое понятие. Термин выделен \term{цветом}. Определения стоит знать
дословно~— на них опирается всё остальное.
\end{opredelenie}

\begin{keyeq*}
\text{Ключевая формула, которую нужно помнить}
\end{keyeq*}

\begin{glavnoe}
Главный вывод фрагмента в одну-две фразы. Если времени мало, читайте
только эти блоки и итоги.
\end{glavnoe}

\begin{smysl}
Объяснение «на пальцах»: что формула означает физически, откуда берётся
тот или иной множитель.
\end{smysl}

\begin{primer*}[Числовой пример]
Подстановка реальных чисел. Помогает почувствовать порядок величин.
\end{primer*}

\begin{vnimanie}
Типичная ошибка или тонкое место, где легко запутаться.
\end{vnimanie}

\begin{podrobnee}[Дополнительный материал]
Длинные выводы и детали. Их можно пропустить при первом чтении и вернуться
к ним при подготовке к экзамену.
\end{podrobnee}

\begin{itogi}
Короткий конспект раздела~— удобно повторять перед экзаменом.
\end{itogi}

\begin{proverka}
\item Вопросы для самопроверки. Ответы на вопросы в конце главы.
\end{proverka}

\begin{masshtab}[дислокация]{-9.6/-9/{ядро}, -8/-5/{расстояние между дислокациями}}
Где «живёт» явление: на логарифмической линейке от 0,1~нм до 1~м
подсвечены характерные размеры (здесь~— ядро дислокации и расстояние
между дислокациями). Каждое деление~— в 10 раз больше предыдущего.
\end{masshtab}

\paragraph{Сноски «Исправление».} В исходном учебнике встречаются опечатки
в формулах и числах. Там, где мы их исправили, стоит сноска с пометкой
\textcolor{fmRed}{\sffamily\bfseries Исправление}~— чтобы при параллельном
чтении расхождение не сбивало с толку.

\paragraph{Обозначения.} Постоянная Больцмана обозначается $k$
($k=8{,}617\cdot10^{-5}\eV/\text{К}$; при комнатной температуре
$kT\approx 0{,}025\eV=1/40\eV$). Энергии дефектов: верхний индекс
$\mathrm F$~— образование (\emph{formation}), $\mathrm m$~— миграция
(\emph{migration}); нижний индекс $v$~— вакансия, $i$~— межузельный атом.
Векторы набраны жирным шрифтом: $\vb b$, $\vb t$.

\clearpage
\section*{Лестница масштабов}
\addcontentsline{toc}{section}{Лестница масштабов}

Материаловедение постоянно перескакивает между масштабами: атомная
плоскость, дислокация, зерно, образец. Одна и та же фраза («кристалл
скользит по плоскости $(111)$») относится сразу к нескольким из них, и
если не держать масштаб в голове, картина рассыпается. Ниже собраны
характерные размеры всего, о чём идёт речь в книге. Ось логарифмическая:
каждое деление~— в 10 раз больше предыдущего, так что от атома до
детали~— всего десять шагов. В тексте у ключевых механизмов стоят
плашки «Масштаб» с такой же линейкой.

\begin{center}
% Лестница масштабов (рисуется вручную; используется во frontmatter.tex).
% Ось x — десятичный логарифм длины в метрах: x = 1,3 см · (lg L + 10).
\begin{tikzpicture}[x=1.3cm, y=1cm, font=\sffamily\scriptsize]
  \def\rowh{0.43}
  % ---- оси и сетка
  \foreach \e in {0,...,10}{\draw[atomLine!15] (\e,0.45) -- (\e,-10.6);}
  \foreach \e/\t in {0/{0,1 нм}, 1/{1 нм}, 2/{10 нм}, 3/{100 нм}, 4/{1 мкм}, 5/{10 мкм},
                     6/{100 мкм}, 7/{1 мм}, 8/{1 см}, 9/{10 см}, 10/{1 м}}{
    \node[atomLine, anchor=south] at (\e,0.45) {\t};
    \node[atomLine, anchor=north] at (\e,-10.6) {\t};}
  \draw[atomLine, line width=0.6pt] (0,0.45) -- (10,0.45) (0,-10.6) -- (10,-10.6);
  % ---- строки: \row{y}{lo}{hi}{цвет}{подпись}
  \newcommand{\scalerow}[5]{%
    \pgfmathsetmacro{\lo}{#2+10}\pgfmathsetmacro{\hi}{#3+10}%
    \fill[#4!75, rounded corners=1.5pt] (\lo,{#1-0.13}) rectangle (\hi,{#1+0.13});
    \pgfmathtruncatemacro{\wide}{(\hi-\lo>=4.5) ? 1 : 0}%
    \pgfmathtruncatemacro{\rside}{(\hi>6.5) ? 1 : 0}%
    \ifnum\wide=1 \node[white, font=\sffamily\bfseries\scriptsize] at ({(\lo+\hi)/2},#1) {#5};
    \else\ifnum\rside=1 \node[anchor=east, #4!80!black, inner sep=2pt] at (\lo,#1) {#5};
    \else \node[anchor=west, #4!80!black, inner sep=2pt] at (\hi,#1) {#5};\fi\fi}
  \newcommand{\scalegroup}[3]{%
    \fill[#2!6] (-0.1,{#1+0.32}) rectangle (10.1,{#1-#3*\rowh+0.1});}
  % ---- Строение
  \scalegroup{0}{fmBlue}{4}
  \node[anchor=east, fmBlue!80!black, font=\sffamily\bfseries\footnotesize] at (10,0.08) {СТРОЕНИЕ};
  \scalerow{-0.25}{-10}{-9.6}{fmBlue}{атомы; межатомное расстояние 0,25 нм}
  \scalerow{-0.68}{-9.6}{-9.2}{fmBlue}{элементарная ячейка 0,3--0,6 нм}
  \scalerow{-1.11}{-6}{-3.5}{fmBlue}{зёрна поликристалла 1 мкм -- 0,3 мм}
  \scalerow{-1.54}{-3}{0}{fmBlue}{образцы и детали}
  % ---- Дефекты
  \scalegroup{-2.05}{fmRed}{6}
  \node[anchor=east, fmRed!80!black, font=\sffamily\bfseries\footnotesize] at (10,-1.97) {ДЕФЕКТЫ (гл.~2)};
  \scalerow{-2.3}{-9.7}{-9.2}{fmRed}{вакансия с зоной искажений}
  \scalerow{-2.73}{-9.6}{-9}{fmRed}{ядро дислокации; толщина границы зерна}
  \scalerow{-3.16}{-9}{-8}{fmRed}{дефект упаковки (ширина расщеплённой дислокации)}
  \scalerow{-3.59}{-9}{-6}{fmRed}{выделения второй фазы, зоны Гинье--Престона}
  \scalerow{-4.02}{-8}{-5}{fmRed}{расстояние между дислокациями}
  \scalerow{-4.45}{-6.5}{-5}{fmRed}{дислокационные ячейки, субзёрна}
  % ---- Процессы и свойства
  \scalegroup{-4.95}{fmGreen}{8}
  \node[anchor=east, fmGreen!80!black, font=\sffamily\bfseries\footnotesize] at (10,-4.87) {ПРОЦЕССЫ И СВОЙСТВА (гл.~1--3)};
  \scalerow{-5.2}{-9.45}{-9.25}{fmGreen}{длина волны электрона у уровня Ферми}
  \scalerow{-5.63}{-9}{-8.3}{fmGreen}{пробег фононов при 300\,К}
  \scalerow{-6.06}{-7.55}{-7.25}{fmGreen}{пробег электронов в Cu при 300\,К}
  \scalerow{-6.49}{-7.2}{-6.8}{fmGreen}{стенка Блоха в железе}
  \scalerow{-6.92}{-8}{-5}{fmGreen}{линии и полосы скольжения}
  \scalerow{-7.35}{-7}{-4}{fmGreen}{кристаллы мартенсита}
  \scalerow{-7.78}{-6}{-3}{fmGreen}{магнитные домены}
  \scalerow{-8.21}{-4.3}{-3.3}{fmGreen}{диффузия за 1 ч у $T_\text{пл}$}
  % ---- Наблюдение
  \scalegroup{-8.71}{fmGray}{4}
  \node[anchor=east, fmGray!80!black, font=\sffamily\bfseries\footnotesize] at (10,-8.63) {ЧЕМ СМОТРЕТЬ};
  \scalerow{-8.96}{-10}{-6}{fmGray}{просвечивающий электронный микроскоп}
  \scalerow{-9.39}{-8.5}{-3}{fmGray}{растровый электронный микроскоп}
  \scalerow{-9.82}{-6.3}{-3}{fmGray}{оптический микроскоп}
  \scalerow{-10.25}{-4}{0}{fmGray}{невооружённый глаз}
\end{tikzpicture}

\end{center}

Как пользоваться: прежде чем представлять себе механизм, найдите на
лестнице его строку и спросите себя, \emph{во сколько раз} он больше атома
и меньше зерна. Например, ступенька скольжения в 100~нм~— это примерно
400 межатомных расстояний, но в 100--1000 раз меньше зерна; поэтому на
поверхности зерна она выглядит тонкой линией.

